\section{Signed Certificate Exceptions}

Signed capability certificates are allowed, but they are not the default authority mechanism. They are adapters at the edges of the registry-first model.

\subsection{Appropriate uses}

Signed certificates \SHOULD{} be limited to:

\begin{itemize}
  \item permit-style approvals;
  \item off-chain orders;
  \item relayed transactions;
  \item cross-chain attestations;
  \item oracle signatures;
  \item temporary sessions;
  \item batch authorization.
\end{itemize}

\subsection{Certificate adapter pattern}

A certificate adapter takes a signed object and turns it into registry state or consumes registry state.

Example:

\begin{lstlisting}
(call-contract! tx TokenA permit Alice Dex USDC 1000 signedPermit)
\end{lstlisting}

\code{permit} does:

\begin{enumerate}
  \item verify signature;
  \item verify domain separation;
  \item verify nonce;
  \item create or update registry allowance;
  \item consume nonce;
  \item emit \code{Approval}.
\end{enumerate}

Then future calls use the normal registry path:

\begin{lstlisting}
(call-contract! tx TokenA transferFrom Alice Bob USDC 10)
\end{lstlisting}

\subsection{Certificate requirements}

\req{SCS-CERT-001 --- Exceptional use}{Signed certificates \MAY{} be supported but \SHOULDNOT{} be the normal authorization mechanism.}

\req{SCS-CERT-002 --- Canonical payload}{A signed certificate \MUST{} have a canonical serialization.}

\req{SCS-CERT-003 --- Domain separation}{A certificate \MUST{} bind the signature to chain, runtime, contract, resource, operation class, and relevant nonce domain.}

\req{SCS-CERT-004 --- Issuer authority}{The verifier \MUST{} check that the signer is authorized to issue the certified authority.}

\req{SCS-CERT-005 --- Non-malleability}{Changing holder, resource, rights, limit, expiry, nonce, or delegation policy \MUST{} invalidate the certificate.}

\req{SCS-CERT-006 --- Replay protection}{A state-changing certificate \MUST{} consume a nonce, certificate identifier, sequence number, or equivalent registry record.}

\req{SCS-CERT-007 --- Registry integration}{Where practical, successful certificate verification \SHOULD{} create, update, or consume registry authority.}

\req{SCS-CERT-008 --- Short validity}{Certificate-only authority \SHOULD{} be short-lived.}

\req{SCS-CERT-009 --- No direct state authority}{A signed certificate \MUSTNOT{} by itself bypass method-policy enforcement or guarded storage.}

\subsection{Hybrid mode}

If a signed certificate requires revocation, spend tracking, or nonce tracking, it becomes hybrid:

\begin{lstlisting}
signature proves issuance
registry proves freshness / non-revocation / unused nonce / remaining spend
\end{lstlisting}

This hybrid mode is preferred for most DeFi certificate flows.
